\chapter{Homotopy and Derived Categories}
\label{ch:derived-categories}

The category of complexes remembers all the differentials. Its homotopy
category identifies maps differing by a cochain homotopy. Its derived
category also makes every quasi-isomorphism invertible. We first keep track
of the maps and signs in the homotopy category, so that the long exact
sequence obtained after localization has the ordinary connecting maps.

\section{Shifts, cones, and triangles}

\begin{proposition}[Shifts and the homotopy category]\label{res:K01}
Let \(\mathcal A\) be an additive category. For a complex \(K\) and
\(m\in\mathbb Z\), put
\[
 (K[m])^n=K^{n+m},\qquad d_{K[m]}^n=(-1)^m d_K^{\,n+m}.
\]
The shifts act on the homotopy category \(\mathsf K(\mathcal A)\) and
compose by addition. If \(\mathcal A\) is abelian, the cohomology functors
on complexes factor through \(\mathsf K(\mathcal A)\), and the
identifications \(H^n(K[m])=H^{n+m}(K)\) are natural.
\end{proposition}
\begin{proof}
The shifted differential squares to zero. If \(f-g=d_L s+s d_K\), where
\(s^n:K^n\to L^{n-1}\), the components \((-1)^m s^{n+m}\) give the
corresponding homotopy on the shifts. Morphisms homotopic to zero form an
additive ideal because precomposition or postcomposition of a
null-homotopy by a cochain map gives a null-homotopy for the composite. Quotienting each
Hom group by that ideal defines \(\mathsf K(\mathcal A)\). On cycles,
the term \(s d_K\) is zero and \(d_L s\) factors through boundaries, so
homotopic maps induce the same cohomology map. Multiplying a differential
by \(-1\) leaves its kernel and image subobjects unchanged. This gives
the cohomology identification and its naturality. See
\cite[Tags 013H and 05RS]{stacks-derived} and
\cite[Section 4.1]{yekutieli}.
\end{proof}

\begin{proposition}[The mapping cone and its maps]\label{res:K02}
For a cochain map \(f:X\to Y\) in an additive category, define
\[
 C(f)^n=Y^n\oplus X^{n+1},\qquad
 d_{C(f)}(y,x)=(d_Y y+f x,-d_X x).
\]
There are cochain maps \(i_f:Y\to C(f)\) and \(p_f:C(f)\to X[1]\),
the inclusion and the projection. Given maps \(a:X\to X'\), \(b:Y\to
Y'\), and a homotopy \(bf-f'a=d_{Y'}s+s d_X\), the formula
\[
 (y,x)\longmapsto (b y+s x,a x)
\]
defines a cone map \(C(f)\to C(f')\). It intertwines the inclusions and
projections with \(b\) and \(a[1]\), respectively.
\end{proposition}
\begin{proof}
In the square of the cone differential the only potentially nonzero
off-diagonal entry is \(d_Y f-f d_X\), which vanishes. For the claimed
cone map, equality of the two off-diagonal entries of the chain map
equation is \(d_{Y'}s+f'a=bf-sd_X\). This is precisely the homotopy
equation. The other entries are the equations for \(a\) and \(b\).
The intertwining identities are immediate from the inclusion and
projection matrices. In an exactly commuting square, take \(s=0\).
For exactly commuting squares this construction preserves composition.
For a square specified only up to homotopy, the formula uses the
chosen homotopy. See \cite[Tags 014E, 014F, and 08RI]{stacks-derived}.
\end{proof}

\begin{proposition}[Contractibility and the cone criterion]\label{res:K03}
The cone of the identity map is contractible. In an abelian category,
\(f\) is a quasi-isomorphism if and only if \(C(f)\) is acyclic.
There exist quasi-isomorphisms of module complexes which are not
homotopy equivalences.
\end{proposition}
\begin{proof}
On \(C(1_X)\) the homotopy \(h^n(y,x)=(0,y)\) satisfies
\(d h+h d=1\). For any \(f\) there is a degreewise split short exact
sequence
\[
 0\longrightarrow Y\xrightarrow{i_f} C(f)
       \xrightarrow{p_f} X[1]\longrightarrow 0.
\]
The connecting map \(H^n(X[1])\to H^{n+1}(Y)\) is \(H^{n+1}(f)\).
Indeed, lifting a cycle of \(X[1]\) by the section of the projection
leaves \(f x\) in the first component after applying the differential. The
long exact sequence proves the criterion in both directions.

For the last assertion, put \(\mathbb Z\xrightarrow{2}\mathbb Z\)
in degrees \(-1,0\). The augmentation to \((\mathbb Z/2)[0]\) is a
quasi-isomorphism. Every map back is zero because every homomorphism
\(\mathbb Z/2\to\mathbb Z\) is zero. Its composite with the augmentation
cannot be homotopic to the identity of the stalk since they have
different \(H^0\). This is the module example from
\cite[Chapter 10]{p103}. The cone criterion is also recorded in
\cite[Tag 05RT]{stacks-derived}.
\end{proof}

\begin{proposition}[Degreewise split sequences]\label{res:K04}
Suppose \(0\to A\xrightarrow{u}B\xrightarrow{v}C\to0\) is degreewise
split. Choose graded identifications \(B^n=A^n\oplus C^n\) compatible
with \(u,v\), and write
\[
 d_B(a,c)=(d_A a+\alpha c,d_C c).
\]
Then \(\alpha:C\to A[1]\) is a cochain map. Its homotopy class is
independent of the splitting and natural in maps of short exact sequences.
The cone row \((u,i_u,-p_u)\) is canonically isomorphic in
\(\mathsf K(\mathcal A)\) to \((u,v,\alpha)\).
\end{proposition}
\begin{proof}
The off-diagonal entry of \(d_B^2=0\) gives
\(d_A\alpha+\alpha d_C=0\), with the sign required for \(A[1]\).
If one graded section differs from another by \(u t\), their
off-diagonal terms differ by \(d_A t-t d_C\), which is a
null-homotopy for maps into \(A[1]\). For a map of short exact sequences
\((a,b,c)\), the difference of the two graded lifts \(b s-s'c\) of \(c\)
is \(u' t\). The equation
\(a\alpha-\alpha'c=d_{A'}t-t d_C\) proves naturality in the homotopy
category.

The map \(q:C(u)\to C\) is \(v\) on \(B\) and zero on \(A[1]\).
In the graded coordinates \((a,c,a')\) of \(C(u)\), a chain section
is \(c\mapsto(0,c,-\alpha c)\). The kernel is \(C(1_A)\), which is
contractible by Proposition~\ref{res:K03}, hence \(q\) is a homotopy
equivalence. Under the chain section, \(-p_u\) is precisely \(\alpha\).
The ensuing sign convention is that the third arrow of a cone triangle
is \(-p_f\). It makes the boundary from a split sequence agree with
the usual cocycle-lift boundary. See
\cite[Tags 05SS and 014L]{stacks-derived}.
\end{proof}

\paragraph{Comparison of five exact terms.}
In a commuting diagram of two exact rows \(A\to B\to C\to D\to E\),
the middle vertical map is an isomorphism when the four others
are. This comparison works in an abelian category using the epic
pullbacks described in Chapter~\ref{ch:filtrations}. For injectivity,
let an element of \(C\) map to zero in the bottom row. Its image in
\(D\) is zero by the comparison at \(D\), so it lifts locally to
\(B\). The image of this lift in the bottom row comes from the
bottom \(A\) by exactness. Lift it to the top \(A\); injectivity
at \(B\) shows that the original lift already comes from \(A\),
so its image in \(C\) is zero. For surjectivity, take an element
of the bottom \(C\) and lift its image in the bottom \(D\) to the
top \(D\). That lift dies in the top \(E\) by injectivity at \(E\).
It therefore lifts locally to the top \(C\). The difference from
the desired bottom element comes from the bottom \(B\), and can
be corrected by lifting through the isomorphism at \(B\).
The zero and image conditions descend along the epic covers.

\begin{theorem}[The triangulated structure on the homotopy category]
\label{res:K05}
The shift \([1]\) and the triangles isomorphic to
\[
 X\xrightarrow{f}Y\xrightarrow{i_f}C(f)
     \xrightarrow{-p_f}X[1]
\]
make \(\mathsf K(\mathcal A)\) a triangulated category.
\end{theorem}
\begin{proof}
Every map is the first arrow of its cone triangle. The identity has a
contractible cone, so its row is isomorphic to \(X\to X\to0\to X[1]\).
The class of triangles is defined up to isomorphism.

\emph{Rotation.} The cone of \(i_f\) has coordinates
\((y,x,y')\in Y^n\oplus X^{n+1}\oplus Y^{n+1}\) and differential
\[
 (d_Yy+fx+y',-d_Xx,-d_Yy').
\]
The map \((y,x,y')\mapsto-x\) to \(X[1]\) has chain section
\(x\mapsto(0,-x,fx)\) and contractible kernel \(C(1_Y)\). After this
identification, the last arrow of the cone triangle for \(i_f\) is
\(-f[1]\). This proves rotation with the stated sign. For inverse
rotation, take the cone of \(p_f[-1]:C(f)[-1]\to X\). Its coordinates
are \((x,y,x')\in X^n\oplus Y^n\oplus X^{n+1}\) with differential
\((d_Xx+x',d_Yy+fx',-d_Xx')\). The map to \(Y\) is
\((x,y,x')\mapsto fx-y\), with chain section \(y\mapsto(0,-y,0)\)
and kernel \(C(1_X)\). The inclusion of \(X\) corresponds to \(f\),
and the last map on the section corresponds to \(i_f\). Thus inverse
rotation also preserves the class.

\emph{Morphisms of triangles.} A commutative square in the homotopy
category can be represented by a square of cochain maps together with
a homotopy between the two composites. Proposition~\ref{res:K02}
supplies the third map of triangles.

If the first two components of a morphism of cone triangles are
isomorphisms in \(\mathsf K(\mathcal A)\), its third component is an
isomorphism as well. To see this without using an octahedron, fix a
complex \(T\). If a cochain map \(u:T\to Y\) has \(i_fu\) null-homotopic,
write the homotopy into \(C(f)^{n-1}=Y^{n-1}\oplus X^n\) as \((s,t)\).
Its second component says that \(t:T\to X\) is a cochain map. Its
first component gives \(u-ft=d_Ys+s d_T\). Conversely, \(i_ff\)
is null-homotopic by \(x\mapsto(0,x)\). This proves exactness at
\(\operatorname{Hom}(T,Y)\). Applying the same calculation to the
rotations gives exact rows of five abelian groups with middle group
\(\operatorname{Hom}(T,C(f))\) and adjacent groups from \(X,Y\) and
their shifts. The morphism of triangles induces a map of these rows,
and the four outer vertical maps are isomorphisms. The preceding
five-term comparison makes the middle map an isomorphism for each \(T\). If \(c:C\to C'\)
is the third component, surjectivity for \(T=C'\) gives a right
inverse to \(c\), and injectivity for \(T=C\) makes it a left
inverse. In particular the morphism axiom, applied to identities
on the first two terms, compares any two distinguished triangles
on a fixed first arrow by an isomorphism fixing those terms.

\emph{The octahedron.} For \(X\xrightarrow{f}Y\xrightarrow{g}Z\)
the two maps
\[
 a:C(f)\to C(gf),\quad (y,x)\mapsto(gy,x),\qquad
 b:C(gf)\to C(g),\quad (z,x)\mapsto(z,fx)
\]
extend the squares on the first two terms. They give the first two
arrows of the fourth triangle. In fact, in the coordinates
\((z,x,y',x')\) of \(C(a)\), its differential is
\[
 (d_Zz+gfx+gy',\,-d_Xx+x',\,-d_Yy'-fx',\,d_Xx').
\]
The chain map \(C(a)\to C(g)\), \((z,x,y',x')\mapsto(z,fx+y')\),
has chain section \((z,y')\mapsto(z,0,y',0)\). Its kernel is
\(C(1_{X[1]})\), with coordinates \((x,x')\) and differential
\((-d_Xx+x',d_Xx')\). Consequently the cone triangle for \(a\) is
isomorphic to
\[
 C(f)\xrightarrow{a}C(gf)\xrightarrow{b}C(g)
   \xrightarrow{i_f[1](-p_g)}C(f)[1].
\]
The outer compatibilities are \(a i_f=i_{gf}g\),
\(b i_{gf}=i_g\), \(p_{gf}a=p_f\), and \(p_g b=f[1]p_{gf}\),
all read directly from the coordinate maps. The closing arrow of the
fourth triangle is \((z,y')\mapsto(-y',0)\), so it agrees with the
composite required by the same choice of cone sign. This gives the
octahedron for composable representatives. For prescribed triangles
on \(f,g,gf\), choose the three comparison isomorphisms just constructed
and conjugate the maps between their third terms by those
isomorphisms. These comparison maps fix \(X,Y,Z\) and their shifts,
so all commuting faces remain commuting. The fourth triangle remains
distinguished because it is isomorphic to the one for the explicit
cone maps. Compare \cite[Tags 014C, 014Q, and 014S]{stacks-derived}.
\end{proof}

\begin{proposition}[Exact sequences from triangles]\label{res:K06}
For a distinguished triangle, cohomology gives a natural long exact
sequence if \(\mathcal A\) is abelian. Each of
\(\operatorname{Hom}_{\mathsf K(\mathcal A)}(T,-)\) and
\(\operatorname{Hom}_{\mathsf K(\mathcal A)}(-,T)\) gives an exact
sequence extended by shifts, with arrows reversed in the second case.
\end{proposition}
\begin{proof}
By isomorphism and rotation it suffices to check exactness at the
middle term of a cone row. Cohomology has the asserted exactness by
the split sequence in Proposition~\ref{res:K03}. Rotation appends the
shifted terms. Use the arrows of the triangle, including its
closing sign. They give natural maps under morphisms of triangles.

For the covariant Hom sequence, suppose \(u:T\to Y\) has \(i_fu\)
null-homotopic, with homotopy \((s,t)\) into
\(Y^{n-1}\oplus X^n=C(f)^{n-1}\). Its second component equation says
that \(t:T\to X\) is a cochain map. The first says
\(u-ft=d_Ys+s d_T\), so \([u]\) is in the image of \([f]\).
Conversely \(i_ff\) is null-homotopic by the homotopy \(x\mapsto(0,x)\).
For the contravariant sequence let \(v:Y\to T\) satisfy
\(vf=d_Ts+s d_X\). The map \(C(f)\to T\),
\((y,x)\mapsto vy+sx\), is a cochain map and restricts to \(v\).
The converse again follows from \(i_ff\simeq0\). Naturality follows
from composition. The shifted long sequences follow by applying
these computations to the rotated cone rows.
Compare \cite[Tags 0149 and 05RS]{stacks-derived} and
\cite[Propositions 5.2.6--5.2.7]{yekutieli}.
\end{proof}

\section{Quasi-isomorphisms and localization}

For the remainder of the chapter \(\mathcal A\) is abelian. The
localization is taken in a universe in which the roofs under
consideration form sets. Bounds on complexes, when imposed, are bounds
on their terms unless cohomology is explicitly specified.

\begin{proposition}[The acyclic subcategory]\label{res:D01}
The acyclic complexes form a thick triangulated subcategory of
\(\mathsf K(\mathcal A)\). Its cone class, the maps whose cones belong
to this subcategory, is the class of quasi-isomorphisms.
\end{proposition}
\begin{proof}
Acyclicity means \(H^n=0\) for all \(n\). It is preserved under
isomorphism in the homotopy category and under shifts by
Proposition~\ref{res:K01}, under cones between acyclic objects by
Proposition~\ref{res:K06}, and under retracts because the cohomology of
a retract is a retract of the cohomology. These are the required
closure properties. The assertion about cones is
Proposition~\ref{res:K03}. In particular quasi-isomorphisms have
two-out-of-three because their cohomology maps do.
\end{proof}

\begin{theorem}[Roofs and their refinements]\label{res:D02}
The quasi-isomorphisms in \(\mathsf K(\mathcal A)\) admit calculi of
right and left fractions. In particular a morphism in the localization
can be written
\[
 X\xleftarrow{s}T\xrightarrow{f}Y,\qquad s\text{ a quasi-isomorphism}.
\]
Two roofs are equal precisely when they can be refined to a common
roof, with both arrows agreeing in \(\mathsf K(\mathcal A)\). There is
also the dual presentation \(X\to T\xleftarrow{s}Y\).
\end{theorem}
\begin{proof}
Here are the two fraction conditions for roofs of the displayed
orientation. Given \(f:X\to Z\) and a quasi-isomorphism \(s:Y\to Z\),
set \(W=C(f,-s)[-1]\), where \((f,-s):X\oplus Y\to Z\).
The cone projection, shifted back, gives maps \(W\to X\) and \(W\to Y\).
Their composites with \(f\) and \(s\) agree in \(\mathsf K(\mathcal A)\)
by the cone homotopy. The projection to \(X\) is a quasi-isomorphism
because its kernel is \(C(-s)[-1]\) and the sequence is degreewise split.
This gives the lifting square for a common denominator.

For cancellation, suppose \(s:Y\to Z\) is a quasi-isomorphism and
\(sf=sg\) in \(\mathsf K(\mathcal A)\) for \(f,g:X\to Y\).
Exactness of Hom applied to the cone triangle for \(s\) factors \(f-g\)
through a map \(u:X\to C(s)[-1]\). The projection
\(C(u)[-1]\to X\) is a quasi-isomorphism since \(C(s)[-1]\)
is acyclic, and precomposing with it kills \(f-g\). This is the
cancellation condition. The left calculus follows from the same
argument with cones in place of the shifted cones, equivalently by
applying the right calculus to the opposite category with cochain
degrees reversed.

For clarity, these conditions give the roof description as follows.
Identities and composites of denominators are denominators. For two
successive roofs, apply the lifting square to the middle numerator and
the next denominator. Precompose the first denominator and the last
numerator with the two maps of that square. Any two choices of square
can first be brought to a common denominator by a further lifting
square, then made equal where necessary by cancellation. The same
construction compares a composite after either input is refined.
For transitivity of common refinement, the maps to the middle roof
are quasi-isomorphisms by two-out-of-three. Refine those two maps by a
lifting square. To check associativity, write the roofs as
\(X_{i-1}\xleftarrow{s_i}T_i\xrightarrow{f_i}X_i\) for \(i=1,2,3\).
Choose \(W\to T_1,T_2\) and \(P\to T_2,T_3\) to lift \(f_1\)
against \(s_2\) and \(f_2\) against \(s_3\), respectively.
The maps \(W\to T_1\) and \(P\to T_2\) are denominators.
Next lift \(W\to T_2\) against \(P\to T_2\), obtaining a
denominator \(Q\to W\) and a map \(Q\to P\).
Then \(Q\to W\to T_1\to X_0\) and \(Q\to P\to T_3\to X_3\)
give the same composite roof for both bracketings, by the three
commuting lifting squares. The image of a
roof under a functor \(F\) inverting quasi-isomorphisms has to be
\(F(f)F(s)^{-1}\). That formula is unchanged under refinement and
respects composition. This proves the localization property and
the claimed equality criterion. See also
\cite[Tags 05R2, 05R6, and 05RU]{stacks-derived} and
\cite[Sections 6.1--6.2]{yekutieli}.
\end{proof}

\begin{theorem}[The derived category]\label{res:D03}
Let \(\mathsf D(\mathcal A)\) be \(\mathsf K(\mathcal A)\) localized at
quasi-isomorphisms. It is the Verdier quotient by the acyclic
subcategory. The quotient \(Q\) preserves distinguished triangles,
and the \(H^n\) factor through \(Q\). A morphism in \(\mathsf D(\mathcal A)\)
is invertible exactly when all its cohomology maps are invertible.
\end{theorem}
\begin{proof}
Write \(W\) for the quasi-isomorphisms. They are the maps whose cones
are acyclic, by Proposition~\ref{res:D01}. Both fraction calculi
hold by Theorem~\ref{res:D02}. This also gives an additive
localization: for two roofs with the same ends take a common
denominator and add the numerators. Any two such choices have a
common further refinement, which proves independence and the
abelian group identities. When precomposing two roofs written
over the same denominator, use the same lifting square for both.
When postcomposing, lift the two numerators against the next
denominator over a common refinement. The sum of their lifts
is a lift of the sum. This proves bilinearity and makes \(Q\)
additive. The zero object stays zero: in a roof out of zero the
denominator has acyclic source and is refined by the map from zero,
and in a roof to zero the numerator is already zero. The biproduct
equations therefore pass to the localization.

For each integer \(m\), the shift \([m]\) preserves \(W\) by
Proposition~\ref{res:K01}. Shifting both arrows of a roof thus
gives the descended shifts, with the same addition law on
shift indices and with the canonical commutation isomorphisms
for \(Q\). Consider now a morphism between distinguished triangles
in \(\mathsf K(\mathcal A)\) whose first two components belong to
\(W\). In each degree the five consecutive groups or objects
\[
 H^n(X)\longrightarrow H^n(Y)\longrightarrow H^n(Z)
       \longrightarrow H^{n+1}(X)\longrightarrow H^{n+1}(Y)
\]
form an exact row by Proposition~\ref{res:K06}, where \(X,Y,Z\)
are the first three objects of the triangle. Naturality gives a
map of exact rows. Both maps on the left and both maps on the
right are isomorphisms, so the five-term comparison just proved
makes the middle map an isomorphism. The third component therefore belongs to \(W\).
In particular a commuting square on the first arrows with vertical
maps in \(W\) extends by Theorem~\ref{res:K05} to a morphism of
triangles all of whose components lie in \(W\).

These give \texttt{IsCompatibleWithTriangulation} for \(W\) in the
fixed Mathlib localization theorem. The required left calculus is
Theorem~\ref{res:D02}. The zero objects and preadditive structures
are those described above. Shifts are additive before localization
by the termwise addition of maps, and afterward because shifting
the sum of roofs over a common denominator gives the sum of their
shifts over the shifted denominator. The commutation of \(Q\) with
these shifts was constructed above. These are the inputs to
\cite[\texttt{CategoryTheory.Triangulated.Localization}]{mathlib-triangulated}.
Its declarations \texttt{pretriangulated},
\texttt{isTriangulated\_functor}, and \texttt{isTriangulated}
give the pretriangulated structure with distinguished triangles
isomorphic to images of distinguished triangles, exactness of
\(Q\), and the octahedral axiom, respectively. The starting
category satisfies the triangulated axioms by Theorem~\ref{res:K05}.

To verify the Verdier universal property, an exact functor
inverting \(W\) kills each acyclic object, since \(0\to C\) lies
in \(W\) for acyclic \(C\). Conversely let an exact functor \(F\)
kill all acyclic objects. For \(s\in W\) it sends the cone
triangle of \(s\) to a distinguished triangle whose third
object is zero. The first arrow of such a triangle is an
isomorphism. Indeed, representable Hom is exact by rotation and
the morphism axiom. For the latter, if \(A\to B\to C\) is the start
of a distinguished triangle and a map \(T\to B\) composes to zero
in \(C\), extend the square from \(T\to0\) to \(B\to C\) to a
map from \(T\to0\to T[1]\xrightarrow{-1}T[1]\) to the rotated
triangle. The commuting last square, shifted back, gives the
factorization through \(A\). Successive arrow composites vanish
by comparing an identity triangle to the given triangle and then
rotating. With third term zero, the resulting Hom sequence makes
the first arrow an isomorphism on every representable Hom, hence
an isomorphism. Consequently \(F\) inverts \(W\).

By Theorem~\ref{res:D02}, \(F\) factors uniquely through \(Q\), by
sending a roof \((s,f)\) to \(F(f)F(s)^{-1}\). The descended
functor is additive by the common-denominator sum. The natural
commutation with shifts of \(F\) is natural on roofs as well:
the squares for \(F(s)\) may be inverted, then composed with the
squares for \(F(f)\). It takes a triangle isomorphic to an image
under \(Q\) to a triangle isomorphic to the corresponding image
under \(F\). Thus the factorization is exact. Its shift data and
the factorizations of natural transformations are determined by
their restrictions to \(Q\): every object is represented in
\(\mathsf K(\mathcal A)\) and naturality on inverses of denominators
follows from naturality on denominators. This proves the exact
universal property for the acyclic subcategory.

Each \(H^n\) inverts quasi-isomorphisms by definition. The localization
property in Theorem~\ref{res:D02} makes its factorization unique.
If a roof represents a map \(w\) with all \(H^n(w)\) invertible, then
the numerator has all \(H^n\) invertible because the denominator does.
Both arrows of the roof therefore become isomorphisms, so \(w\) is an
isomorphism. The converse is functorial. In particular an object of
the derived category with all cohomology zero is zero.
The long cohomology sequence from Proposition~\ref{res:K06} descends
to derived triangles by their description as localized cone rows.
\end{proof}

\begin{theorem}[Triangles from short exact sequences]\label{res:D04}
Every short exact sequence \(0\to A\xrightarrow{u}B\xrightarrow{v}C
\to0\) of complexes induces a distinguished triangle
\[
 A\xrightarrow{u}B\xrightarrow{v}C
   \xrightarrow{\delta}A[1]\quad\text{in }\mathsf D(\mathcal A).
\]
The construction is natural in maps of short exact sequences and its
cohomology sequence has the usual connecting homomorphisms.
\end{theorem}
\begin{proof}
The map \(q:C(u)\to C\), \(q(b,a)=v(b)\), is a chain epimorphism
whose kernel is \(C(1_A)\). Thus \(Q(q)\) is invertible. Transporting
the cone triangle defines
\(\delta=Q(-p_u)Q(q)^{-1}\). A map of short exact sequences induces
the strict cone map from Proposition~\ref{res:K02} with \(s=0\).
It intertwines \(q\) and \(-p_u\). Inverting \(Q(q)\) proves
naturality without any choice of a splitting. A cone cocycle
\((b,a)\) projecting to the given cohomology class satisfies
\(d_Bb=-u(a)\). The ordinary boundary is the class of \(-a\),
which is induced by \(-p_u\). In an abelian category this is the
same diagram chase in the pullback of \(Z^n(C)\) along \(B^n\to
C^n\). The resulting boundary kills boundaries and is given by
\(-p_u\) after the isomorphism on cohomology.
See \cite[Tag 0152]{stacks-derived}.
\end{proof}

\section{Good truncation and bounded models}

\begin{proposition}[Good truncation]\label{res:D05}
For any cochain complex \(K\), the good truncations are
\[
\begin{aligned}
 (\tau_{\le n}K)^i&=
  \begin{cases}
   K^i & i<n,\\
   \ker(d_K^n) & i=n,\\
   0 & i>n,
  \end{cases}
&
 (\tau_{\ge n}K)^i&=
  \begin{cases}
   0 & i<n,\\
   \operatorname{coker}(d_K^{n-1}) & i=n,\\
   K^i & i>n.
  \end{cases}
\end{aligned}
\]
They preserve cohomology on the indicated side and give zero on
the other side. They induce functors on the derived category and a
natural triangle
\[
 \tau_{\le n}K\longrightarrow K\longrightarrow\tau_{\ge n+1}K
       \longrightarrow(\tau_{\le n}K)[1].
\]
Shifts commute with these constructions, with the truncation indices
translated by the shift.
\end{proposition}
\begin{proof}
The incoming differential at the kernel factors through it by
\(d^2=0\). The outgoing differential at the cokernel descends for
the same reason. The edge cohomology groups are both the relevant
cycle object modulo the preceding boundaries. The other cohomology
degrees follow from the displayed complexes. Chain maps between
the truncations inherit the induced kernel or cokernel map. For the
upper truncation a homotopy restricts at \(n\) to the source cycles,
where the discarded term \(s^{n+1}d^n\) vanishes. For the lower
truncation the degree \(n+1\) homotopy component is followed by the
projection onto the target cokernel in degree \(n\). This projection
kills the discarded term \(d^{n-1}s^n\) in degree \(n\). These observations make truncation
functorial on the homotopy category. The cohomology computation
shows it preserves quasi-isomorphisms, so it descends to
\(\mathsf D(\mathcal A)\).

The quotient \(K/\tau_{\le n}K\) maps naturally to
\(\tau_{\ge n+1}K\). Its kernel is the two-term identity complex on
\(\operatorname{im}d_K^n\) in degrees \(n,n+1\), since
\(K^n/\ker(d_K^n)=\operatorname{im}d_K^n\). Thus the quotient map
is a quasi-isomorphism, and Theorem~\ref{res:D04} supplies the
triangle. Naturality on complexes passes to the localization because
the triangle maps commute with quasi-isomorphisms and therefore
with their inverses. For example
\((\tau_{\le n}K)[m]\cong\tau_{\le n-m}(K[m])\), since multiplying
the differential by a sign does not change its kernel. The same
calculation applies to the lower truncation.
See \cite[Tags 05RV and 08J5]{stacks-derived} and
\cite[Definition 7.3.6]{yekutieli}.
\end{proof}

\begin{proposition}[Strict models and bounded localizations]\label{res:D06}
Write \(\mathsf K^+,\mathsf K^-,\mathsf K^b\) for the homotopy
categories with the indicated bounds on terms, and
\(\mathsf D^+,\mathsf D^-,\mathsf D^b\) for the full subcategories
with the corresponding bounds on cohomology. Localization of each
\(\mathsf K^\ast\) at its quasi-isomorphisms gives an equivalence
with \(\mathsf D^\ast\), for \(\ast\in\{+,-,b\}\). These are
triangulated subcategories of the corresponding unbounded
categories.
\end{proposition}
\begin{proof}
Upper and lower cohomological bounds are represented by strict ones
by Proposition~\ref{res:D05}. If both bounds are imposed, the two
truncations can be applied consecutively. Their bounds are preserved
under shifts and any cone of a map between objects with the chosen
type of bounds again has those bounds (for cohomological bounds,
use the long exact sequence). Hence the full subcategories are
triangulated. The fraction constructions of Theorem~\ref{res:D02}
use only shifts, cones and finite sums, and thus give fraction
calculi on the corresponding bounded homotopy categories.

For fullness, take a right roof between strict bounded representatives
\(X,Y\), with middle complex \(T\). If \(X,Y\) are zero below \(a\),
then \(H^i(T)=0\) for \(i<a\) because the denominator is a
quasi-isomorphism. Both maps from \(T\) factor through
\(T\to\tau_{\ge a}T\), since their degree \(a\) maps kill
\(\operatorname{im}d_T^{a-1}\). That quotient is a quasi-isomorphism.
If \(X,Y\) are zero above \(b\), precompose both maps by
\(\tau_{\le b}T\to T\), a quasi-isomorphism by the upper bound on
\(H^i(T)\). With both bounds apply both operations. For faithfulness,
an equality of two right roofs has a common refinement in
\(\mathsf K\). The strict bounded complexes that occur in the two
roofs admit a common bound of the relevant type. Truncating the
refining complex exactly as above keeps the comparison maps and
the quasi-isomorphisms. Equality of composites in the homotopy
category is preserved for upper truncation by precomposition.
For lower truncation a homotopy into a complex zero below \(a\)
descends, because its degree \(a\) component has zero target and
its degree \(a+1\) component has unchanged source. Thus the
truncated refinement proves equality in the bounded localization.
See \cite[Tags 05RQ, 05RV, and 05RW]{stacks-derived}.
\end{proof}

\begin{theorem}[The standard heart]\label{res:D07}
The cohomological subcategories
\[
 \mathsf D^{\le0}=\{K:H^i(K)=0\text{ for }i>0\},\qquad
 \mathsf D^{\ge0}=\{K:H^i(K)=0\text{ for }i<0\}
\]
form the standard \(t\)-structure. Its heart, their intersection, is
equivalent to \(\mathcal A\) by the stalk embedding \(A\mapsto A[0]\)
and the inverse \(H^0\).
\end{theorem}
\begin{proof}
Proposition~\ref{res:K01} gives
\(\mathsf D^{\le0}[1]\subseteq\mathsf D^{\le0}\) and
\(\mathsf D^{\ge0}[-1]\subseteq\mathsf D^{\ge0}\). For
orthogonality take \(X\in\mathsf D^{\le0}\), \(Y\in\mathsf D^{\ge1}\).
Represent \(X,Y\) by complexes supported in degrees at most \(0\)
and at least \(1\), respectively. In a right roof \(X\leftarrow T
\to Y\), replace \(T\) by \(\tau_{\le0}T\) without changing the
roof, since \(X\leftarrow T\) is a quasi-isomorphism. The map from
this new middle term to \(Y\) is zero degree by degree. This proves
\(\operatorname{Hom}(X,Y)=0\). The truncation triangle for \(n=0\)
provides the remaining \(t\)-structure axiom.

If an object has cohomology only in degree zero, apply the two good
truncations at zero to obtain its stalk \(H^0(K)[0]\) with canonical
quasi-isomorphisms in the truncation zigzag. Between two stalks
the roof computation of Proposition~\ref{res:D06} can be made
with both bounds zero. Its denominator is then an ordinary
isomorphism in \(\mathcal A\). Thus the stalk embedding is fully
faithful and \(H^0\) is its inverse on the heart. In particular for
stalks \(A,B\) and \(m<0\) we have
\(\operatorname{Hom}_{\mathsf D}(A[0],B[m])=0\) by orthogonality.
See \cite[Tags 05RV and 06XS]{stacks-derived} and
\cite[Proposition 7.5.1]{yekutieli}.
\end{proof}

\begin{proposition}[The bounded Postnikov tower and splitting]
\label{res:D08}
An object \(K\) with bounded cohomology is obtained from its cohomology
stalks by the finite tower of triangles
\[
 \tau_{\le i-1}K\longrightarrow\tau_{\le i}K
 \longrightarrow H^i(K)[-i]\xrightarrow{\epsilon_i}
            (\tau_{\le i-1}K)[1].
\]
If
\(\operatorname{Ext}^{\,i-j+1}_{\mathcal A}(H^i(K),H^j(K))=0\)
for all \(j<i\), then \(K\) is isomorphic to
\(\bigoplus_i H^i(K)[-i]\). Here Ext denotes the shifted Hom groups
between stalks. The comparison with resolutions is given in
Theorem~\ref{res:R11}. The splitting obtained from these
vanishing conditions need not be canonical.
\end{proposition}
\begin{proof}
Apply Proposition~\ref{res:D05} to \(\tau_{\le i}K\) and cut at
\(i-1\). The upper and lower pieces are the stated truncation and
the degree \(i\) cohomology stalk. The tower begins at zero and
ends at \(K\) because of the cohomological bounds. The maps
\(\epsilon_i\) record the attachments.

Inductively suppose \(\tau_{\le i-1}K\cong
\bigoplus_{j<i} H^j(K)[-j]\). The components of \(\epsilon_i\) then
belong to the groups
\[
 \operatorname{Hom}_{\mathsf D(\mathcal A)}
       \bigl(H^i(K)[0],H^j(K)[\,i-j+1\,]\bigr).
\]
The direct sum is finite and all these groups vanish by assumption.
Thus \(\epsilon_i=0\). Applying
\(\operatorname{Hom}_{\mathsf D}(H^i(K)[-i],-)\) to the triangle
gives a section of the map to \(H^i(K)[-i]\). Indeed the Hom
exactness in Proposition~\ref{res:K06} holds in any triangulated
category by the triangle axioms. If \(g u=0\) for \(g\) the
second arrow, the square from the zero triangle
\(T\to 0\to T[1]\xrightarrow{-1}T[1]\) to the rotated target
triangle lifts by the morphism axiom, and unshifting the lift
expresses \(u\) through the first arrow. The contravariant case
is dual. The section and the inclusion from the preceding
truncation now give a map from their direct sum to
\(\tau_{\le i}K\). It is an isomorphism on cohomology by the long
exact sequence, so it is an isomorphism by
Theorem~\ref{res:D03}. The section is a choice, which proves the
claim without assigning a preferred splitting.
Compare \cite[Tags 08J5, 0GM4, and 0EWX]{stacks-derived}.
\end{proof}
